% Section 7 -- Related Work
% Target: 2 pages.

DLC sits at the intersection of three research lineages: authorization
logic, capability-based access control, and AI agent delegation.

\subsection{Authorization logics with \texttt{says}-modality}

Garg \& Pfenning's constructive authorization logic~\cite{garg2006csf}
provides DLC's intuitionistic foundation: \texttt{says} as a modality
in a sequent calculus with cut elimination and non-interference proven
on paper.  Nexus Authorization Logic (NAL)~\cite{schneider2011nal}
mechanizes a related says-modality system in $\sim$3000 LOC of Coq.
DLC extends this lineage with explicit cryptographic carrier
semantics (\textsc{says-I} carries the signature intrinsically),
linear obligations ($\Box_O$), time modality ($\Diamond_\tau$), and
IFC labels.  Where NAL proves soundness against a belief-style
semantics, DLC's T2 targets a cryptographic correspondence---
arguably a stronger semantic anchor.

\subsection{Proof-carrying authorization}

Appel \& Felten~\cite{appel1999pcfs} introduce the
\emph{proof-carrying authorization} (PCA) paradigm: each request is
accompanied by a logical proof of its authorization.  PCFS, Bauer's
PCA-on-the-web work, and Garg's proof-carrying file system extend
this.  All use higher-order logic that is not decidable.  DLC's
contribution is occupying the decidability sweet-spot: the
propositional fragment admits a mechanized structural checker (T1),
and the design targets an $O(|M| \cdot \log |\Gamma|)$ bound for the
full calculus --- a target that is stated, not proven, at this
revision (Section~\ref{sec:metatheory}).

\subsection{Authenticated delegation for AI agents}

The agent-economy line begins with South et al.~(ICML
2025)~\cite{south2025authdelegation}: authenticated delegation as
OAuth+OIDC extensions.  AITH (Chen,
2026)~\cite{chen2026aith} mechanizes five Tamarin theorems
including post-quantum security.  LLMbda Calculus (Garby et al.,
2026)~\cite{garby2026llmbda} provides termination-insensitive
non-interference for an LLM-invoking lambda calculus.  DLC's
positioning: AITH covers the symbolic-protocol axis; LLMbda covers
the type-theoretic axis; DLC unifies both, plus the wire format and
cryptographic correspondence, into a single artifact with the
same chain-splicing-defeating rule appearing at each layer.

\subsection{Capability tokens}

Macaroons~\cite{birgisson2014macaroons} and
biscuits~\cite{couprie2022biscuits} provide cryptographically
verifiable, attenuable capability tokens.  ZCAP-LD~\cite{w3c2021zcap}
and UCAN~\cite{ucan2024} extend the capability-token model to
linked-data contexts.  These systems excel at wire-format design;
they lack a typed logic of obligations or non-interference.  DLC's
\texttt{attenuate} rule with the IFC-lattice side condition is the
natural lift of macaroon caveats to a typed setting.

\subsection{OAuth and the RFC 8693 lineage}

RFC 8693 (Token Exchange)~\cite{rfc8693} introduces the nested
\texttt{act} claim for delegation chains.  The February--March 2026 IETF
OAuth mailing-list thread on chain splicing~\cite{ietf2026splicing}
documented the cross-validation gap: an STS validating
\texttt{subject\_token} and \texttt{actor\_token} independently can
issue tokens asserting chains that never occurred.
draft-ietf-oauth-attenuating-agent-tokens-00~\cite{niyikiza2026attn}
proposes mitigations.  DLC's \texttt{delegate} rule is the
structural answer the IETF mitigation patterns implement procedurally:
the same-principal condition is a type-system invariant, not a
runtime check.

\subsection{Mechanized cryptography in Lean}

VCVio (Tuma et al.,
2026)~\cite{tuma2026vcvio} mechanizes EUF-CMA security for Schnorr
signatures in Lean 4 using monadic oracle effects.  This is the
closest Lean-side precedent for T2's computational reduction.  Earlier
Lean cryptographic work (\texttt{SymbolicCryptographyLean}) and
the Mathlib4 categorical infrastructure provide the substrate.

\subsection{Symbolic protocol verification with multiple provers}

Cremers \& Sasse~\cite{cremers2013tamarin} provide Tamarin.
Blanchet~\cite{blanchet2022proverif} provides ProVerif.  Joint use
of both has been done case-by-case (e.g., for TLS 1.3), but the
\emph{methodological argument} that running new protocols through
two provers detects model-overfitting bugs is, to our knowledge, not
made systematically in the delegation literature.  This paper
contributes one such systematic case: the two-modeling-gaps story
in Section~\ref{sec:symbolic}.
